\section{Artinian rings}
\label{section-artinian}

\noindent
Artinian rings, and especially local Artinian rings,
play an important role in algebraic geometry, for example
in deformation theory.

\begin{definition}
\label{definition-artinian}
A ring $R$ is {\it Artinian} if it satisfies the
descending chain condition for ideals.
\end{definition}

\begin{lemma}
\label{lemma-finite-dimensional-algebra}
Suppose $R$ is a finite dimensional algebra over a field.
Then $R$ is Artinian.
\end{lemma}

\begin{proof}
The descending chain condition for ideals obviously holds.
\end{proof}

\begin{lemma}
\label{lemma-artinian-finite-nr-max}
If $R$ is Artinian then $R$ has only finitely many maximal ideals.
\end{lemma}

\begin{proof}
Suppose that $\mathfrak m_i$, $i = 1, 2, 3, \ldots$ are
pairwise distinct maximal ideals.
Then $\mathfrak m_1 \supset \mathfrak m_1\cap \mathfrak m_2
\supset \mathfrak m_1 \cap \mathfrak m_2 \cap \mathfrak m_3 \supset \ldots$
is an infinite descending sequence (because by the Chinese
remainder theorem all the maps $R \to \oplus_{i = 1}^n R/\mathfrak m_i$
are surjective).
\end{proof}

\begin{lemma}
\label{lemma-artinian-radical-nilpotent}
Let $R$ be Artinian. The Jacobson radical of $R$ is a nilpotent ideal.
\end{lemma}

\begin{proof}
Let $I \subset R$ be the Jacobson radical.
Note that $I \supset I^2 \supset I^3 \supset \ldots$ is a descending
sequence. Thus $I^n = I^{n + 1}$ for some $n$.
Set $J = \{ x\in R \mid xI^n = 0\}$. We have to show $J = R$.
If not, choose an ideal $J' \not = J$, $J \subset J'$ minimal (possible
by the Artinian property). Then $J'/J$ is a simple $R$-module, hence
isomorphic to $R/\mathfrak m$ for some maximal ideal $\mathfrak m$,
see Lemma \ref{lemma-characterize-length-1}.
Then $\mathfrak m I^n$ kills $J'$. Since $I \subset \mathfrak m$
we conclude that $I^{n + 1} = I^n$ kills $J'$. Hence $J' = J$
which is a contradiction.
\end{proof}

\begin{lemma}
\label{lemma-product-local}
Any ring with finitely many maximal ideals and
locally nilpotent Jacobson radical is the product of its localizations
at its maximal ideals. Also, all primes are maximal.
\end{lemma}

\begin{proof}
Let $R$ be a ring with finitely many maximal ideals
$\mathfrak m_1, \ldots, \mathfrak m_n$.
Let $I = \bigcap_{i = 1}^n \mathfrak m_i$
be the Jacobson radical of $R$. Assume $I$ is locally nilpotent.
Let $\mathfrak p$ be a prime ideal of $R$.
Since every prime contains every nilpotent
element of $R$ we see
$ \mathfrak p \supset \mathfrak m_1 \cap \ldots \cap \mathfrak m_n$.
Since $\mathfrak m_1 \cap \ldots \cap \mathfrak m_n \supset
\mathfrak m_1 \ldots \mathfrak m_n$
we conclude $\mathfrak p \supset \mathfrak m_1 \ldots \mathfrak m_n$.
Hence $\mathfrak p \supset \mathfrak m_i$ for some $i$, and so
$\mathfrak p = \mathfrak m_i$. Thus the spectrum of $R$
is the discrete topological space $\{\mathfrak m_1, \ldots, \mathfrak m_n\}$.
By Lemma \ref{lemma-disjoint-implies-product} applied $n - 1$ times
we find that $R = R_1 \times \ldots \times R_n$ where the
spectrum of $R_i$ is a singleton for each $i$.
Thus $R_i$ is a local ring and since it is a localization of $R$
(by the lemma), it is one of the local rings of $R$ as desired.
\end{proof}

\begin{lemma}
\label{lemma-artinian-finite-length}
A ring $R$ is Artinian if and only if it has finite length
as a module over itself. Any such ring $R$ is both Artinian and
Noetherian, any prime ideal of $R$ is a maximal ideal, and $R$ is equal
to the (finite) product of its localizations at its maximal ideals.
\end{lemma}

\begin{proof}
If $R$ has finite length over itself then it satisfies both
the ascending chain condition and the descending chain
condition for ideals. Hence it is both Noetherian and Artinian.
Any Artinian ring is equal to product of its localizations
at maximal ideals by Lemmas \ref{lemma-artinian-finite-nr-max},
\ref{lemma-artinian-radical-nilpotent}, and \ref{lemma-product-local}.

\medskip\noindent
Suppose that $R$ is Artinian. We will show $R$ has finite
length over itself. It suffices to exhibit a chain of
submodules whose successive quotients have finite length.
By what we said above
we may assume that $R$ is local, with maximal ideal $\mathfrak m$.
By Lemma \ref{lemma-artinian-radical-nilpotent} we have
$\mathfrak m^n =0$ for some $n$.
Consider the sequence
$0 = \mathfrak m^n \subset \mathfrak m^{n-1} \subset
\ldots \subset \mathfrak m \subset R$. By Lemma
\ref{lemma-dimension-is-length} the length of each subquotient
$\mathfrak m^j/\mathfrak m^{j + 1}$ is the dimension of this
as a vector space over $\kappa(\mathfrak m)$. This has to be
finite since otherwise we would have an infinite descending
chain of sub vector spaces which would correspond to an
infinite descending chain of ideals in $R$.
\end{proof}









